\section{Getting started}
\label{sec:start}

\subsection{Your first typeset}

Do this now, with the document you are reading.

\begin{enumerate}[leftmargin=*]
  \item Click \ui{Open project} in the toolbar and choose the folder
        containing this file (\path{demo-project/le-guide}).
  \item In the \ui{Files} panel, click \texttt{main.tex}.
  \item Press \kbd{Cmd+Enter}, or click \ui{Typeset}.
\end{enumerate}

The PDF appears on the right. The \ui{Issues} panel along the bottom fills with
anything the run reported.

\subsection{You will need to press it twice}
\label{sec:passes}

The first build of this guide produces a document whose table of contents is
empty and whose cross-references read \texttt{??}. Press \ui{Typeset} again and
both fill in.

That is not a bug, and it is not specific to this editor --- it is how \LaTeX{}
works. Cross-references, the table of contents and page numbers are written to
an auxiliary file on one pass and read back on the next. Tools like
\texttt{latexmk} hide this by re-running the compiler until the output settles.
LaTeXEditor runs \emph{one} pass per press, which is faster and makes the error
output easier to read, but leaves the second press to you.

Two practical consequences:

\begin{itemize}[leftmargin=*]
  \item \textbf{Cross-references and the ToC need two presses} after they
        change. Just press \ui{Typeset} again.
  \item \textbf{BibTeX and Biber are not run.} A \verb|\bibliography{refs}|
        will never resolve, because the editor does not invoke a second
        program. This guide therefore uses a \texttt{thebibliography}
        environment written directly into \texttt{main.tex}, which resolves in
        a single pass. If your project needs full BibTeX, run it once in a
        terminal, then carry on typesetting here.
\end{itemize}

If you would rather have the loop handled for you, select \texttt{latexmk} in
the compiler menu --- it performs its own re-runs internally.

\subsection{Where the output goes}

Everything the compiler emits --- \texttt{.pdf}, \texttt{.aux}, \texttt{.log},
\texttt{.toc} --- is written to a \texttt{build/} folder beside your project,
never next to your sources. Your project folder stays readable.

\subsection{Which file gets typeset}

Not, in general, the one you are looking at. Before each run the editor scans
the project for the \texttt{.tex} file containing \verb|\documentclass| and
typesets \emph{that}.

This is what you want in a multi-file project. Editing
\texttt{sections/03-writing.tex} and pressing \ui{Typeset} builds
\texttt{main.tex}, which pulls the section back in via \verb|\input|. You never
have to navigate back to the root file to build.

\subsection{Saving}

There is no save button and you rarely need \kbd{Cmd+S}. The editor writes your
file one second after you stop typing. A dot beside the filename in the toolbar
means there are changes not yet on disk; it disappears when they are. Switching
files flushes any pending write immediately, so edits made in the last second
before you click away are not lost.

Turn on \ui{Auto} in the toolbar and every save also triggers a typeset --- a
tight write-and-see loop, at the cost of a compiler run every time you pause.

\subsection{Starting from a template}

\ui{New from template} offers a gallery: Basic Article, IEEE Conference Paper,
ACM Paper, Thesis/Dissertation, Beamer Presentation, CV/Resume, Cover Letter and
Homework/Problem Set. Choosing one writes a complete, compiling project into
your current folder and opens its main file. It is a faster start
than an empty buffer, and a reasonable way to see how a document class you have
not used before is normally set up.
